\section{Physical consistency, literature comparison and readiness}
\subsection{What the literature does and does not establish here}
\paragraph{Ideal Weyssenhoff.} Trautman \cite{Trautman} supplies the minimal variational setting; the specific ideal effective source/transport is selected from \cite{BHL,OK}. The latter is a constitutive field formulation, not a proof of a causal radial collapse problem for any EOS. Our spherical cancellation, complete reduced characteristic pencil, endpoint Jordan analysis, general-EOS derivative and contortion-jump test are derived consequences checked independently. They are not claimed as novel physics or a new fundamental no-bounce theorem.

\paragraph{Published collapse comparisons.} Hashemi et al. \cite{Hashemi} underlie Model A's homogeneous averaged closure. Pop{\l}awski \cite{Pop} uses a spherical thermodynamic setting with additional choices, including particle production in the discussed scenarios. A synchronous comoving Tolman gauge restricts acceleration and hence pressure gradients; it does not reproduce our arbitrary $\Phi(t,r)$ matter problem. Luz--Mena--Ziaie \cite{LMZ} study an \emph{effective dust} condition: ordinary pressure cancels spin pressure. That is different from ordinary $p=0$, $P=-q$. Neither their conclusions nor a static Weyssenhoff stellar solution in \cite{LC} supplies the missing radial hyperbolicity or exterior for our EOS.

\paragraph{Geometric well-posedness is not matter closure.} The recent analysis of Luz--Mena \cite{LWP} proves a geometric EC Cauchy result in a generalized harmonic gauge with compatible, sufficiently regular matter supplied ab initio; the stated hypotheses include conservation constraints on those supplied sources. Its non-strictly hyperbolic Leray--Ohya/Gevrey analysis is not a proof that our EOS/spin evolution yields those regular sources or a strongly hyperbolic numerical matter system. It does not contradict the explicit nonreal acoustic roots in Eq.~\eqref{eq:charpoly}. We do not claim EC geometry lacks a local well-posedness theory.

\paragraph{Modern spin hydrodynamics.} Weickgenannt \cite{Weick} derives a kinetic, dissipative spin theory and obtains causal, linearly stable modes about homogeneous equilibrium under specified matching and transport-coefficient conditions. Xie et al. \cite{Xie} analyze another dissipative spin framework and show that necessary small/large-wave-number conditions can miss finite-wave-number instabilities. These are distinct constitutive systems, not conflicting theorems about every first-order spin theory. They demonstrate the need for spin thermodynamics, microscopic matching and transport data, not a license to add relaxation/viscosity to ideal Weyssenhoff matter. We neither import their coefficients into EC nor establish their validity in dense gravitational collapse.

\subsection{Readiness decision for each formulation}
\begin{longtable}{p{.26\textwidth}p{.29\textwidth}p{.34\textwidth}}
\toprule Formulation & What is established & Decision\\\midrule\endhead
Model A: averaged GSM & Conditional local action, source, variance transport and characteristic pencil derived. & Algebraically closed under explicit averaging assumptions; same dust/enthalpy obstruction and unresolved finite boundary. The FLRW reference stays unchanged.\\
GR dust/perfect fluid & Standard spinless limits and analytical benchmarks. & Reference equations retained; a selected gauge/center/boundary problem precedes any new prototype.\\
Model B: coherent radial ideal fluid, $0<w<1$, $z<z_s$ & Formal closed continuum system and real causal material eigenbasis away from endpoints. & Candidate local smooth patch only; full coupled initial-boundary formulation and microscopic EOS still require justification. No prototype authorized by all six gates.\\
Both: ordinary dust, $0<z<1/2$ & Formal closure but imaginary acoustic pair. & Generic smooth radial evolution blocked.\\
Both: $h=0$ or $z>z_h$, $w<1$ & Singular time matrix or superluminal radial branch. & No causal bounce-through-domain solver justified.\\
Stiff ideal branch $w=1$ & Luminal pair for $h\ne0$; ratio fixed along shells. & Separate physical/energy/boundary problem; no generic regular-bounce claim.\\
Common-congruence statistical second moments & Conditional transport closes extra advected moments; same acoustic pair. & Not a stabilization mechanism. Physical averaging assumptions must be justified.\\
General spin/velocity/geometry correlations & Symmetry permits additional spherical sources. & Missing constitutive moments and evolution; no closed principal symbol or solver.\\\bottomrule
\end{longtable}
The task's six implementation prerequisites have \emph{not} all been met for a finite high-density full-Weyssenhoff collapse problem. We therefore stop before a spherical EC numerical prototype. This is the requested scientific result, not an unfinished implementation. A favorable local acoustic patch alone is insufficient to meet the gauge, EOS, center and exterior gates.

\subsection{What a justified simulation would record}
For any future accepted formulation, record $R$ and neighboring material shell labels, $U$, $M$, $\Gamma$, $B$, $\Phi$, $\rho,p,b,q$, all equation residuals and shell proper time $\tau_r=\int A\,\dd t$. Areal extrema require $R>0$, $U=0$ and the proper sign of $\Dt U$; a volume minimum instead requires $\Theta=0$, $\Dt\Theta>0$. A regular bounce additionally needs nonsingular matter/torsion/curvature, a continuing well-posed solution, finite particle Jacobians and compatible boundary data. An isolated turning point is not that result.

Track $\sigma^2$, $\Theta$, material caustics, noncentral $R\to0$, metric and full-connection curvature, and separate numerical/chart failure from physical divergence. Recollapse and oscillations require resolved extrema sequences with refinement and valid equations throughout. Geometric diagnostics remain
\begin{equation}
 \theta_\pm=\frac2R(U\pm\Gamma),\quad
 \theta_+\theta_-=\frac4{R^2}(2M/R-1),\quad
 N_R=\ln(R/R_i),\quad N_V=\tfrac13\ln(BR^2/B_iR_i^2).
\end{equation}
Both negative expansions identify future trapped spheres, both positive identify anti-trapping. A horizon-regular chart is needed for an implemented exterior; $2M/R=1$ is not itself a curvature singularity. Define accelerated volume e-folds, if used, by integrating $\Theta/3$ only when $\Theta>0$ and $\Dt\Theta/3+\Theta^2/9>0$, using each shell's proper time. Center values use regular limits. Export metric and spin-junction residuals; until a matching construction exists the status is \emph{not computed}. None of these local diagnostics establishes an event horizon, eternal oscillations or a baby universe.

\subsection{Open physics and the Dirac research direction}
The precise open tasks are: independently reconstruct the constrained full matter action; identify a microscopic high-density EOS and polarization susceptibility; justify a coarse-graining scale and close any retained spin/velocity/metric moments; establish the coupled gauge and characteristic problem; derive regular material surfaces or a justified non-vacuum exterior; and specify any surface stress/spin action without undefined products of distributions. These are separate obligations; minimal EC gravity alone does not solve them.

A minimally coupled Einstein--Cartan--Dirac action is a justified \emph{next research comparison}, since it specifies spin current and spin dynamics microscopically instead of guessing a fluid completion. Its spin source is a totally antisymmetric axial current, generally not the convective $u^\lambda b_{\mu\nu}$ source here. Eliminating its algebraic torsion produces a local axial-current contact interaction and a nonlinear Dirac equation. That interaction is algebraic in the spinors and does not change the Dirac principal symbol $\gamma^\mu\xi_\mu$; its characteristic cone remains metric-null. This structural observation is not a derivation of all gravitational constraints or a nonlinear existence theorem in a chosen spherical gauge.

Spherical spinors require careful symmetry of bilinears and sources; a single classical spinor is not a rotationally invariant scalar field. Torsion-free Einstein--Dirac work such as \cite{FSY} uses paired spin states to achieve spherical stress. In an EC extension, the \emph{total} axial current and cross-species interactions must be computed before imposing a singlet/multiplet ansatz. If it cancels identically for classical fields, the corresponding classical torsion effect cancels too. An unpolarized quantum state with zero mean current but nonzero current correlations is not equivalent to setting the classical current to zero before eliminating torsion. Quantum expectation values then require a controlled state, renormalization and correlation closure. Thus EC--Dirac offers a more definite starting action, not a guaranteed bounce or an automatic shortcut around the statistical-matter problem. No Dirac solver or invented four-fermion coefficient is implemented in this stage.

\subsection{Automated validation and its limits}
The new tests check the full reduced characteristic pencils, eigenvectors and endpoint Jordan structure, EOS derivative, sign domains and stiff exception; a nonspherical shear/acceleration control; direct radial divergence cancellation; spin transport and conserved particle polarization; center-series mass and lapse constraints; the full contortion-jump Riemann coefficient and its vanishing Ricci contraction; mixed Einstein components and the conditional geodesic bounce bound. Stage 1's 23 checks are rerun for geometry, Bianchi, mass/ADM constraints, conservation, Frenkel/source coefficients, FLRW, Schwarzschild and LTB limits, as are the existing application tests.

These are algebraic, analytical and numerical-evaluation validations, not a PDE time evolution, proof of microscopic validity or global well-posedness. Exact test inventory and environment versions are recorded with reproduction commands. Stage 1 files and the FLRW runtime remain intact; this directory is a new research stage.
